\documentclass{article}
\usepackage{mathnotes}

\title{Set Theory}

\begin{document}

\section{Set Theory}

\begin{definition}[Set]\synonyms{collection, family}
A \textbf{set} is a collection of objects, considered as a whole.
\end{definition}

\begin{definition}[Element]\synonyms{member, point}
The objects that make up a \dref{set} are called its \textbf{elements} or its \textbf{members.}
\end{definition}

\begin{definition}[Membership criterion]
The \textbf{membership criterion} for a \dref{set} $X$ is a statement of the form $x \in X \iff P(x),$ where $P(x)$ is a proposition that is true for precisely for those objects $x$ that are \dref{elements} of $X,$ and no others.
\end{definition}

\begin{axiom}[Axiom of Extensionality]\label{axiom-of-extensionality}
Two \dref{sets} are equal if and only if they have the same \dref{elements}. Formally, for any sets $A$ and $B$:

\[ A = B \iff (\forall x)(x \in A \iff x \in B) \]
\end{axiom}

\begin{definition}[subset]
If $X$ and $Y$ are \dref{sets} such that every element of $X$ is also an element of $Y,$ then we say $X$ is a \textbf{subset} of $Y,$ denoted as $X \subset Y.$ Formally,

\[ X \subset Y \iff (\forall x)(x \in X \implies x \in Y) \]
\end{definition}

\begin{definition}[superset]
If $X$ and $Y$ are \dref{sets} such that every element of $Y$ is also an element of $X,$ then we say $X$ is a \textbf{superset} of $Y,$ denoted as $X \supset Y.$ This is the same as $Y \subset X.$
\end{definition}

\begin{theorem}\label{set-equality-via-subset-inclusion}
Two sets $X$ and $Y$ are equal if and only if $X$ is a subset of $Y$ and $Y$ is a subset of $X.$
\end{theorem}

\begin{proof}
Suppose $X$ and $Y$ are sets with $X \subset Y$ and $Y \subset X.$ Now, suppose $x \in X.$ Then, $x \in Y.$ Conversely, suppose $y \in Y.$ Then $y \in X.$ Thus, $(\forall x)(x \in X \iff x \in Y),$ and \dref[$X = Y$]{axiom-of-extensionality}.
\end{proof}

\begin{definition}[Function]\synonyms{mapping}
Consider two \dref{sets}, $A$ and $B,$ whose \dref{elements} may be any objects whatsoever, and suppose that with each \dref{element} $x$ of $A$ there is associated, in some manner, any \dref{element} of $B,$ which we denote by $f(x).$ Then $f$ is said to be a \textbf{function} from $A$ to $B.$
\end{definition}

\begin{definition}[Domain of Definition]
If $f$ is a \dref{function} from the \dref{set} $A$ to the \dref{set} $B,$ the \dref{set} $A$ is called the \textbf{domain} of $f.$
\end{definition}

\begin{definition}[Value]
If $f$ is a \dref{function} from the \dref{set} $A$ to the \dref{set} $B,$ the \dref{elements} $f(x) \in B$ are called the \textbf{values} of $f.$
\end{definition}

\begin{definition}[Range]
The \dref{set} of all \dref{values} of a \dref{function} $f$ is called the \textbf{range} of $f.$
\end{definition}

\begin{definition}[Sequence]
A \textbf{sequence} is a \dref{function} $f$ defined on the \dref{set} $J$ of all positive integers.
\end{definition}

\begin{note}\label{sequence-notation}
If $f(n) = x_n,$ for $n \in J,$ it is customary to denote the \dref{sequence} $f$ by the symbol $\{x_n\},$ or sometimes by $x_1, x_2, x_3, \dots.$
\end{note}

\begin{definition}[Term]
If $f$ is a \dref{sequence} denoted as $\{x_n\},$ the \dref{values} of $f,$ that is, the \dref{elements} $x_n,$ are called the \textbf{terms} of the \dref{sequence}.
\end{definition}

\begin{note}\label{sequence-terms-not-distinct}
The \dref{terms} of a \dref{sequence} need not be distinct.
\end{note}

\begin{definition}[Inverse Image]
If $f: A \to B$ and $E \subset(B),$ then $f^{-1}(E)$ denotes the \dref{set} of all $x \in A$ such that $f(x) \in E.$ We call $f^{-1}(E)$ the \textbf{inverse image} of $E$ under $f.$
\end{definition}

\subsection{De Morgan's Laws}

\begin{theorem}\label{complement-of-union-is-intersection-of-complements}
The complement of a union is equal to the intersection of complements.
\end{theorem}

\begin{proof}
Let $A$ and $B$ be sets. We want to show that

\[ (A \cup B)^c = A^c \cap B^c. \]

Suppose $x \in (A \cup B)^c.$ Then, if $x \in A$ or $x \in B,$ then $x \in A \cup B$ and $x \notin (A \cup B)^c,$ a contradiction. Therefore, $x \notin A$ and $x \notin B.$ That is, $x \in A^c$ and $x \in B^c,$ therefore $x \in A^c \cap B^c.$
\end{proof}

\begin{theorem}\label{complement-of-intersection-is-union-of-complements}
The complement of an intersection is equal to the union of complements.
\end{theorem}

\begin{proof}
Let $A$ and $B$ be sets. We want to show that

\[ (A \cap B)^c = A^c \cup B^c. \]

Suppose $x \in (A \cap B)^c.$ Then, $x$ is not in $A \cap B,$ that is, $x$ is either not in $A$ or it is not in $B$ or it is in neither. If $x$ is not in $A,$ then it is in $A^c,$ and therefore it is in $A^c \cup B^c.$ The same approach works with $B,$ and therefore $x \in  A^c \cup B^c,$ and we have shown $(A \cap B)^c = A^c \cup B^c.$
\end{proof}

\end{document}
